% ==============================================================================
% فایل تنظیمات و بسته‌های نگارش پایان‌نامه (Preamble)
% دانشگاه تهران - دانشکدگان فنی - دانشکده مهندسی برق و کامپیوتر
% طراحی سیستم تأمین داده جهت استنتاج WCET
% دانشجو: محمد مهدی دوست محمدی (810100142)
% ==============================================================================

% تنظیمات ابعاد صفحه و حاشیه‌ها (مطابق فایل الگوی گزارش PSA1 و دانشگاه تهران)
\usepackage[
    a4paper,
    top=30mm,
    bottom=25.4mm,
    left=22.2mm,
    right=24.3mm,
    headheight=16pt,
    footskip=12mm
]{geometry}

% تنظیم عمق فهرست مطالب و شماره‌گذاری بخش‌ها (تنها فصول، بخش‌ها و زیربخش‌ها)
\setcounter{secnumdepth}{2}
\setcounter{tocdepth}{2}

% بسته‌های ریاضی و نمادها
\usepackage{amsmath}
\usepackage{amssymb}
\usepackage{amsfonts}
\usepackage{mathtools}

% بسته‌های درج تصاویر، نمودارها و رنگ‌ها
\usepackage{graphicx}
\usepackage{xcolor}
\usepackage{tikz}
\usetikzlibrary{shapes.geometric, arrows.meta, positioning, calc, fit, backgrounds}
\usepackage{eso-pic}

% بسته‌های جداول پیشرفته
\usepackage{booktabs}
\usepackage{tabularx}
\usepackage{multirow}
\usepackage{array}
\usepackage{float}
\usepackage{caption}
\usepackage{subcaption}

% بسته‌های نمایش کد و الگوریتم
\usepackage{listings}
\usepackage{tcolorbox}
\tcbuselibrary{skins, breakable}

% تنظیمات سربرگ و پانویس (Header & Footer)
\usepackage{fancyhdr}
\pagestyle{fancy}
\fancyhf{}
\fancyhead[R]{\small\slshape \leftmark}
\fancyhead[L]{\small\thepage}
\fancyfoot[C]{}
\renewcommand{\headrulewidth}{0.6pt}
\renewcommand{\footrulewidth}{0pt}

% استایل صفحات ابتدایی فصل‌ها
\fancypagestyle{plain}{
    \fancyhf{}
    \fancyfoot[C]{\small\thepage}
    \renewcommand{\headrulewidth}{0pt}
}

% تعریف کادربندی ظریف دور صفحه (مطابق با کادربندی موجود در PSA1_FP_Report)
\AddToShipoutPictureBG{%
  \begin{tikzpicture}[remember picture, overlay]
    \draw[line width=0.6pt, color=gray!80]
      ($(current page.north west) + (1.2cm, -1.2cm)$)
      rectangle
      ($(current page.south east) + (-1.2cm, 1.2cm)$);
  \end{tikzpicture}%
}

% تنظیمات لینک‌ها و بوکمارک‌های پی‌دی‌اف
\usepackage{hyperref}
\hypersetup{
    pdfpagelabels=false,
    colorlinks=true,
    linkcolor=blue!75!black,
    citecolor=green!60!black,
    urlcolor=blue!85!black,
    filecolor=magenta,
    pdfstartview={FitH},
    pdfauthor={محمد مهدی دوست محمدی},
    pdftitle={طراحی سیستم تأمین داده جهت استنتاج WCET},
    pdfsubject={پایان‌نامه کارشناسی مهندسی برق - الکترونیک},
    pdfkeywords={WCET, ARM Cortex-M7, ETM, CoreSight, OpenCSD, STM32H7, Instruction Trace}
}

% تنظیمات نمایش کدهای برنامه‌نویسی (C, Python, Assembly)
\definecolor{codegreen}{rgb}{0,0.6,0}
\definecolor{codegray}{rgb}{0.5,0.5,0.5}
\definecolor{codepurple}{rgb}{0.58,0,0.82}
\definecolor{backcolour}{rgb}{0.97,0.97,0.97}

\lstdefinestyle{customc}{
    backgroundcolor=\color{backcolour},
    commentstyle=\color{codegreen},
    keywordstyle=\color{blue}\bfseries,
    numberstyle=\tiny\color{codegray},
    stringstyle=\color{codepurple},
    basicstyle=\ttfamily\footnotesize,
    breakatwhitespace=false,
    breaklines=true,
    captionpos=b,
    keepspaces=true,
    numbers=left,
    numbersep=5pt,
    showspaces=false,
    showstringspaces=false,
    showtabs=false,
    tabsize=2,
    frame=single,
    rulecolor=\color{black!20}
}
\lstset{style=customc}

% سازگاری کرنل جدید لاتک با ماژول‌های bidi در TeX Live 2026
\providecommand{\UseMathForPositioningText}{}
\providecommand{\UseTaggingSocket}[1]{}

% نادیده‌گرفتن نویسه نامرئی U+200F جهت سازگاری با فونت‌های تک‌فاصله
\usepackage{newunicodechar}
\newunicodechar{‏}{}

% ------------------------------------------------------------------------------
% بسته زی‌پرشین (باید همواره آخرین بسته فراخوانی‌شده باشد)
% ------------------------------------------------------------------------------
\usepackage{xepersian}

% تنظیم قلم‌های استاندارد زبان فارسی و انگلیسی (دو سایز بزرگتر و خواناتر)
\settextfont[
    Path=fonts/,
    BoldFont=BNaznnBd.ttf,
    Scale=1.22
]{BNazanin.ttf}
\setlatintextfont[Scale=1.1]{Times New Roman}

% تعریف قلم‌های سفارشی تیتر و عناوین
\defpersianfont\titr[
    Path=fonts/,
    Scale=1.28
]{BTitrBd.ttf}

\defpersianfont\zar[
    Path=fonts/,
    BoldFont=BZarBd.ttf,
    Scale=1.22
]{BZar.ttf}

% افزایش عمومی فاصله بین خطوط در سرتاسر متن پایان‌نامه
\linespread{1.35}

% ارقام فارسی در محیط‌های متنی و ریاضی
\SepMark{.}

% نمادهای بالت استاندارد با قلم ریاضی (جلوگیری از گلیف ناموجود در قلم متن فارسی)
\renewcommand{\labelitemi}{\ensuremath{\bullet}}
\renewcommand{\labelitemii}{\ensuremath{\circ}}
\renewcommand{\labelitemiii}{\ensuremath{\ast}}
\renewcommand{\labelitemiv}{\ensuremath{\cdot}}

% اصلاح عنوان مراجع برای جلوگیری از فونت ناموجود در سربرگ‌های لاتین
\renewcommand{\bibname}{مراجع}
